\section{Source reconciliation, verification, and further research}
\label{audit:section}

\subsection{Precise changes to the research status}

The targeted audit has not established an additional false theorem in
the inspected source dependencies. The proposed corrections concern
the scope and status of research questions, the distinction between
inherited and new results, and the conventions needed when combining
formulas. This distinction matters in a rapidly developing collection:
the date on an archive does not by itself determine which other results
its research list has incorporated.

\paragraph{The reconstruction questions.}
Question R1 in the latest \emph{Quartic Spectral Reconstruction}
report asks for the complete smaller-order-ten family and then the
smaller-order-twelve family \cite{tail-quartic}. Those two specified
tasks are now resolved by Theorems~\ref{t10:complete}
and~\ref{t12:complete}. In particular, the older excluded progression
for order ten is no longer an unresolved family: the inherited
eventual theorem and congruence theorem reduce the remaining task
to the single case $b=962$, which is certified here. For order twelve,
the complete finite range and the proved eventual range cover every
$b\ge7$. The unrestricted assertion for every $a<b$, requested in
R2 of the quartic report, remains open. So does the full fixed-gap-six
family $b=a+3$. The newly included cases $(p,q)=(10,16)$ and $(12,18)$
do not settle that family. Every rank statement is in the specified
formal coordinate space; no arithmetic independence of zeta values
follows from it.

\paragraph{The alternating-tail question.}
The question headed ``Bilinear adjacent-tail certificates'' in the
latest alternating-tail report \cite{tail-alternating} is answered
within exactly its two-product rational ansatz. Theorem~\ref{bi:classification}
gives necessity and sufficiency, and Theorem~\ref{bi:sum} supplies the
finite and infinite endpoint identities. The independently shifted
quadratic test case is therefore completely evaluated, including
confluence, mixed argument derivatives, and anchored rectangular
primitives. The divisor obstruction is an impossibility result for
this certificate class. It is not an impossibility theorem for all
summation identities involving an incompatible pair of denominators.
Larger support, different auxiliary terms, and independently variable
Hurwitz orders remain separate problems.

\paragraph{The Stieltjes question.}
Question 9 of the earlier \emph{Polygamma Tail Identities} report
excluded zero argument-derivative orders from its positive-order
product generator \cite{tail-polygamma}. Theorem~\ref{edge:theorem}
fills the edge with exactly one zero derivative order. Together with
the inherited positive-positive formula, the uniform dictionary
\eqref{edge:uniform-dictionary} covers every same-point quadratic
product with positive total argument-derivative order. All spectral
derivatives in that dictionary are evaluated at the single integer
$w=r+t+1$; ordinary zeta values at other positive integers occur in
the coefficients. The double-undifferentiated germ $H_{1,1}$ and the
nonzero periodic-translation formulas were already available in the
unified manuscript \cite{manuscript}. Equation~\eqref{edge:corner}
checks compatibility with the former. The translated products and
their contact corrections retain their original attribution.

The already incorporated Stieltjes source corrections, general
digamma-power identities, Gamma normal forms, and convolution towers
are not relabelled as new discoveries. The Gaussian $S_6$ and revised
$S_8$ conjectures are also unaffected. Evaluating another conductor-four
tail does not evaluate either of those separately defined periods.

\subsection{Normalization and endpoint requirements for integration}

Three conventions are essential when moving these results into the
unified manuscript. First, the finite summation identities retain their
upper endpoint until its limit is justified. In particular,
\eqref{bi:finite} and \eqref{sh:primitive-finite} record information
that can survive analytic continuation beyond the ordinary convergence
region. A continued expression is not thereby the sum of the
pointwise-specialized summands.

Second, meromorphic expressions must be combined before taking a
removable value or an order derivative. The Mellin factors involving
$s\zeta(s+1)$ are entire combinations, including at $s=0$.
At the Stieltjes edge, both resonant local terms are subtracted
jointly. Substituting a parameter equal to one into a formula proved
only for larger integer orders does not prove the extension.
Likewise, merging two separated singular points changes the local
subtraction problem. It requires a collision theorem, rather than
literal substitution into a translated finite-part formula.

Third, the finite part uses the original coordinate $x$ and the
constant term of the cutoff integral. Integration by parts therefore
has the explicit endpoint law \eqref{edge:ibp}. For example,
\[
 \FP\int_0^1\bigl\{
   \gamma_1(x)\gamma_0'(x)+\gamma_0(x)\gamma_1'(x)
   \bigr\}\,dx=-\zeta(2)-\zeta'(2),
\]
which is not zero. The Stieltjes convention
\eqref{eq:stieltjes-convention} fixes every sign in this statement.
Similarly, differentiating the weighted kernel $K(z,w)$ requires
differentiating its factor $2n+z+w$ as well as both tails. The
additional terms in \eqref{sh:mixed-formula} are part of the identity.
These are requirements for correct reuse, not assertions that the
source theorems violated them.

\subsection{What the verification actually establishes}

The order-ten certificate proves a nonzero retained determinant
$60407$ modulo $65537$, and an independent reconstruction gives
$15578$ for the corresponding original determinant modulo the same
prime. The order-twelve certificate covers all $11513$ integers
$7\le b\le11519$; its two independent implementations agree on every
residue. The infinite remaining range is handled by the written
threshold proof. These are exact arithmetic certificates. The
bilinear/kernel checks include $144$ rational cutoff identities,
$72$ rational kernel cutoff identities, and $3125$ finite divisor
strings, in addition to the symbolic cancellation and determinant
checks. Its $21$ numerical diagnostics at $50$-digit precision have
maximum absolute discrepancy below $5.4\cdot10^{-51}$.
The Mellin verifier records $2002$ exact rational or algebraic
checks and $11$ direct integer-order or scaled quadratures. The
Stieltjes verifier records $81$ exact endpoint identities, four
printed-formula checks, and five direct finite-part quadratures,
with one repeated at a second splitting point. Its largest recorded
absolute discrepancy is less than $1.3\cdot10^{-46}$ at $55$-digit
working precision. The quadratures are numerical diagnostics, not
interval certificates, and do not replace the holomorphy arguments
for arbitrary complex orders or all-index derivatives. The supplied
logs distinguish these two kinds of evidence. The combined
\texttt{verify\_all.py} run completed successfully.

\subsection{Twelve concrete questions beyond the completed results}

The following questions supplement the Mellin-specific programme
already given above. They ask for structural identities, finite
reductions, or complete nonvanishing arguments. No claim of global
novelty or of arithmetic independence is built into their formulation.

\begin{enumerate}[leftmargin=*,label=\textbf{R\arabic*.}]

\item \textbf{A uniform reconstruction mechanism.}
Find a determinant identity or a recurrence proving
$\det M_{a,b}\ne0$ for every $1\le a<b$, or determine an exact
exceptional set if one exists. The completed families through smaller
order twelve constrain any proposed formula. At present the highest
Bernoulli-degree determinant explains eventual nonvanishing for fixed
$a$, but does not exclude cancellation by the lower degrees uniformly
in both parameters. A mechanism preserving those lower-degree terms
would resolve R2 of the quartic report, rather than merely enlarging
the range of computations.

\item \textbf{The moving boundary in the gap-six family.}
Derive a bordered-Hankel or comparable exact formula for $M_{a,a+3}$
that retains all lower-index terms and the omitted-binomial
conventions at the boundary. The fixed-order contraction used in the
source does not justify omitting this information for arbitrarily
large $a$ at the fixed separation. A recurrence in $a$
with explicit nonzero initial data would be particularly useful.
This is a sharply defined intermediate problem between the existing
gap-four theorem and unrestricted unequal orders.

\item \textbf{Closed reconstruction coefficients at orders ten and twelve.}
For $a=5,6$, replace an unevaluated matrix inverse in the reconstruction
of a prescribed zeta coordinate by a uniform coefficient formula.
The target is a fixed-size kernel or product expression for the
polynomial weight, including the ordinary-integral subtraction
constants. Such a formula would turn the existence theorem for each
larger order into an immediately usable family of polygamma
integrals. Producing another determinant ratio is only useful if its
minors themselves admit a proved finite reduction.

\item \textbf{Minimal support beyond the adjacent bilinear class.}
Allow a finite collection of products $E_{n+i}F_{n+j}$ with rational
coefficients and require the difference of the certificate to contain
only a prescribed nonzero bilinear term. Determine the smallest
support needed, modulo identities already implied by the tail
recurrences. Does the divisor obstruction of
Proposition~\ref{bi:divisor} persist, or can an incompatible pair such
as generic separately shifted pure powers be accommodated?
A necessary-and-sufficient criterion would distinguish a genuine new
summation mechanism from a rewriting of the two-product certificate.

\item \textbf{Finite ranks for longer adjacent denominator ladders.}
For
\[
 L_d(z)=\sum_{j\ge0}\frac{(-1)^j}{(j+z)_d},\qquad
 \Re z>0,\quad d\ge2,
\]
determine the minimal separated rank of the weighted two-shift kernels
obtained from admissible certificates. Even $d$ admits a pairing of
adjacent linear factors; odd $d$ encounters the divisor obstruction
at generic independent shifts. The rank-two formula for $d=2$
suggests an exact finite determinant hierarchy. The objective is a
minimal rank theorem and explicit confluent determinants for every
admissible $d$, not just a separable expansion whose length has not
been proved optimal.

\item \textbf{Three independently colored tails.}
Classify finite certificates built from neighboring products of three
tails with independent denominator polynomials and colors. State
which quadratic and linear residual tails must survive and give the
complete finite endpoint. The previously developed powers of a single
tail do not settle this independent-data problem. A first substantive
target is a three-shift quadratic family, with a specified reduction
to depth-two sums and one-factor functions; unrestricted depth-one
closure should not be assumed.

\item \textbf{The continued alternating-square law at negative orders.}
Starting from its exact finite endpoint, construct a jointly
holomorphic subtraction formula at each nonpositive tail order and
derive every order derivative there. The already documented
continued value $1/4$ at order zero, despite vanishing pointwise
summands, is the compulsory normalization test
\cite{tail-alternating}. The desired identities should identify the
finite Euler-polynomial contributions and the spectral logarithmic
terms separately. This would give exact negative-order identities
without confusing continuation with an ordinary series.

\item \textbf{Independent nonperiodic Stieltjes shifts.}
Construct an order-generating formula for the ordinary integrals
\[
 \int_0^1\gamma_m^{(r)}(x+a)\gamma_n^{(t)}(x+b)\,dx,
 \qquad \Re a,\Re b>0,
\]
with $a$ and $b$ independent. These are not the inherited periodic
translations $\{x+a\}$. Begin with the antisymmetric first-derivative
combination, since its symmetric companion is an endpoint derivative.
Specify a finite function class or a genuinely new two-parameter
master, and prove how its diagonal and boundary limits recover the
appropriate coincident finite parts. The independent shift difference
is the missing variable in the classical coincident product kernel.

\item \textbf{Cubic Stieltjes products with independent indices.}
Determine an explicit holomorphic germ for
$\FP\int_0^1\gamma_m^{(r)}\gamma_n^{(t)}\gamma_j^{(v)}\,dx$
when $r+t+v>0$. A first nontrivial target is
$\FP\int_0^1\gamma_1(x)^2\gamma_0'(x)\,dx$.
The all-index digamma-power results concern a different restriction
of this problem. Fourier or Mellin analysis suggests that genuine
multiple-zeta order derivatives may be needed. A successful theorem
should identify that alphabet and all local counterterms before
claiming a finite reduction; the single-spectral-argument quadratic
closure does not imply its cubic analogue.

\item \textbf{Universal relations among the quadratic finite parts.}
For fixed total Stieltjes index and fixed positive total argument
derivative, classify all relations implied by the exact coefficient
dictionary. Determine whether exchange symmetry and the endpoint
identity generate all of them, or whether further relations remain.
The cancellation of the highest zeta derivative in
\eqref{edge:top} gives a concrete first test. Seek a uniform generating
basis and dimension in a declared formal coefficient alphabet. This
would explain cancellations systematically while avoiding an
unproved claim about independence of evaluated constants.

\item \textbf{Simultaneous and iterated triple collisions.}
For three periodically translated Stieltjes factors, compare the
constant term obtained by a simultaneous fixed-coordinate collision
with the two possible successive pairwise collisions. Compute any
difference as an explicit local distribution and test whether a
canonical correction makes the resulting composition associative.
The existing quadratic contact and dilation laws
\cite{manuscript} are mandatory compatibility conditions, not a proof
of this three-factor assertion. The question is meaningful only
after the two-variable cutoff path and the retained finite term are
specified; those choices are part of the proposed theorem.

\item \textbf{A finite evaluation of the logarithmic Gamma correction.}
The inherited function
$S(a,b)=\sum_{n\ge0}\arctan^2(b/(n+a))$ already has an ordinary
Gamma-integral representation and exact shift and distribution laws
\cite{tail-alternating}. Find a proved finite evaluation at a
nonzero mass, for example on a rational-shift locus with
$e^{-2\pi b}$ algebraic, in an explicitly stated collection of
polylogarithm values, order derivatives, and Gamma values. The
existing small-mass series and integral remain valid answers at
generic parameters. A finite special-value reduction would be an
additional identity and would clarify which extra functions the
higher spectral defects actually require.

\end{enumerate}

\subsection{Proposed repository integration}

The accompanying integration note proposes separate inserts for the
certificate classification, the shifted kernels and primitives, the
Mellin resonances, the Stieltjes edge, and the two reconstruction
families. It records their exact source dependencies and the status
sentences that should change. The current manuscript should retain
its established collision, convolution, and Gamma-normal-form
chapters, with cross-references to the added edge theorem. The
complete verification records should travel with the relevant
inserts, including the finite-field residue list and its independent
replay. No repository files are modified by this delivery.
